Este projeto de pesquisa propôs implementar e avaliar métodos de inteligência artificial (IA) para o desenvolvimento de ferramentas automatizadas destinadas ao auxílio no diagnóstico de doenças da retina, especialmente a retinopatia diabética, bem como alterações relacionadas a doenças cardiovasculares e metabólicas. A retinopatia diabética (RD) é uma complicação microvascular específica do diabetes mellitus (DM) \cite{wong_diabetic_2016, wang_diabetic_2018}, e uma das principais causas de perda de visão evitável na população adulta \cite{wong_diabetic_2016, wang_diabetic_2018, melo_current_2018, korn_malerbi_feasibility_2022, teo_global_2021}. Envolve um espectro de lesões microvasculares retinianas, mas evidências crescentes apontam também para a degeneração neuronal e processos inflamatórios como eventos precoces em sua fisiopatologia \cite{wong_diabetic_2016, wang_diabetic_2018}. Foram desenvolvidos modelos baseados em clustering, máquinas de vetores de suporte (SVM) e redes neurais convolucionais (CNN) para classificação automática das imagens retinográficas. Esses modelos foram comparados em diferentes contextos, usando tanto imagens coletadas pelo projeto quanto bases públicas. Os resultados indicam que o uso combinado dessas técnicas pode aprimorar significativamente a triagem e o diagnóstico precoce, proporcionando benefícios substanciais ao Sistema Único de Saúde (SUS), especialmente ao reduzir filas e agilizar encaminhamentos especializados. Além disso, o projeto promove avanços na engenharia biomédica e contribui para a formação de recursos humanos altamente capacitados. \\